% GENERATED FILE. Edit canonical lessons, then run npm run generate:books.
% canonical-source-hash: fnv1a64:700698e07e7bc3e9
% canonical-lessons: TA-C347-ataram, TA-C347-nirupi, TA-C347-tarkkam, TA-C347-atippatai, TA-C347-vilakkam

\chapter{Evidence and Logic}
\label{ch:ta-evidence-and-logic}
\hlchaptermodality{347}

\begin{chapteropening}
\textbf{By the end of this chapter:} \emph{I can ask for evidence and say whether an argument holds together.}

Complete the last lesson of chapter 347: I can ask for evidence and say whether an argument holds together.
\end{chapteropening}

\section[\texorpdfstring{\ta{ஆதாரம்} (ātāram)}{ஆதாரம் (ātāram)}]{\ta{ஆதாரம்} (ātāram) — evidence; proof; a source}

\label{lesson:TA-C347-ataram}



\begin{quote}
Before the new one: say the Tamil for a complaint, then the Tamil for an athlete.
\end{quote}

\subsection*{You'll want to know: \ta{ஆதாரம்}}
\textbf{\ta{ஆதாரம்}} — \emph{ātāram} — \textquotedblleft{}evidence; proof; a source\textquotedblright{}.

\begin{quote}
\ta{உங்களிடம்} \ta{இதற்கு} \ta{ஆதாரம்} \ta{இருக்கிறதா}?
\end{quote}

\begin{quote}
\emph{uṅkaḷiṭam itaṟku ātāram irukkiṟatā?}
\end{quote}

\begin{quote}
Do you have evidence for this?
\end{quote}

Plural: \ta{ஆதாரங்கள்}, \emph{ātāraṅkaḷ}.

With -\ta{உடன்}, a -\ta{அம்} noun takes its oblique -\ta{அத்து}-: \ta{ஆதாரத்துடன்}, with evidence.

\begin{cousinweb}
\textbf{\ta{ஆதாரம்}} (\emph{ātāram}) is borrowed from Sanskrit \emph{ādhāra}, \textquotedblleft{}support, substratum\textquotedblright{} (CDIAL 1165). Tamil gives the final -a the noun ending -\ta{அம்}, and \textquotedblleft{}what holds a claim up\textquotedblright{} became \textquotedblleft{}evidence\textquotedblright{}.
\end{cousinweb}

\begin{grammarlens}[title={What you've built}]
One more word for giving reasons and opinions.
\end{grammarlens}

\subsection*{Your turn}
\begin{itemize}
\raggedright
  \item \emph{Say it:} \emph{ātāram}
  \item \emph{Say it:} \emph{ātāram}, once more
  \item \emph{Say it:} \emph{uṅkaḷiṭam itaṟku ātāram irukkiṟatā?} — Do you have evidence for this?
  \item \emph{Recall:} say the Tamil for a complaint, then the Tamil for an athlete, then say \emph{ātāram} again
\end{itemize}

\begin{tcolorbox}[breakable,colback=teal!4,colframe=teal!35!black,title={Before you move on}]
What does \ta{ஆதாரம்} mean? (\textquotedblleft{}Evidence; proof; a source\textquotedblright{}.) Say it once more.
\end{tcolorbox}
\section[\texorpdfstring{\ta{நிரூபி} (nirūpi)}{நிரூபி (nirūpi)}]{\ta{நிரூபி} (nirūpi) — to prove}

\label{lesson:TA-C347-nirupi}



\begin{quote}
Before the new one: say the Tamil for an athlete, then the Tamil for evidence.
\end{quote}

\subsection*{You'll want to know: \ta{நிரூபி}}
\textbf{\ta{நிரூபி}} — \emph{nirūpi} — \textquotedblleft{}to prove\textquotedblright{}.

\begin{quote}
\ta{அவர்} \ta{தன்} \ta{கருத்தை} \ta{ஆதாரத்துடன்} \ta{நிரூபித்தார்}.
\end{quote}

\begin{quote}
\emph{avar taṉ karuttai ātārattuṭaṉ nirūpittār.}
\end{quote}

\begin{quote}
He proved his point with evidence.
\end{quote}

Past (I …): \ta{நிரூபித்தேன்}, \emph{nirūpittēṉ}.

What is proved takes -\ta{ஐ}: \ta{கருத்தை} \ta{நிரூபி}, to prove a point.

\begin{cousinweb}
\textbf{\ta{நிரூபி}} (\emph{nirūpi}) is a verb Tamil built on Sanskrit \emph{nirūpaṇa}, \textquotedblleft{}stating, determining; examination\textquotedblright{} (Monier-Williams), with the verb ending -\ta{இ} that many Sanskrit loans take, as in \ta{சுவாசி}, \textquotedblleft{}to breathe\textquotedblright{}.
\end{cousinweb}

\begin{grammarlens}[title={What you've built}]
One more word for giving reasons and opinions.
\end{grammarlens}

\subsection*{Your turn}
\begin{itemize}
\raggedright
  \item \emph{Say it:} \emph{nirūpi}
  \item \emph{Say it:} \emph{nirūpi}, once more
  \item \emph{Say it:} \emph{avar taṉ karuttai ātārattuṭaṉ nirūpittār.} — He proved his point with evidence.
  \item \emph{Recall:} say the Tamil for an athlete, then the Tamil for evidence, then say \emph{nirūpi} again
\end{itemize}

\begin{tcolorbox}[breakable,colback=teal!4,colframe=teal!35!black,title={Before you move on}]
What does \ta{நிரூபி} mean? (\textquotedblleft{}To prove\textquotedblright{}.) Say it once more.
\end{tcolorbox}
\section[\texorpdfstring{\ta{தர்க்கம்} (tarkkam)}{தர்க்கம் (tarkkam)}]{\ta{தர்க்கம்} (tarkkam) — logic; reasoning}

\label{lesson:TA-C347-tarkkam}



\begin{quote}
Before the new one: say the Tamil for evidence, then the Tamil for to prove.
\end{quote}

\subsection*{You'll want to know: \ta{தர்க்கம்}}
\textbf{\ta{தர்க்கம்}} — \emph{tarkkam} — \textquotedblleft{}logic; reasoning\textquotedblright{}.

\begin{quote}
\ta{அவன்} \ta{சொன்னதில்} \ta{தர்க்கம்} \ta{இல்லை}.
\end{quote}

\begin{quote}
\emph{avaṉ coṉṉatil tarkkam illai.}
\end{quote}

\begin{quote}
There is no logic in what he said.
\end{quote}

\ta{தர்க்கமாக}, logically: an -\ta{அம்} noun + -\ta{ஆக} makes the adverb.

\begin{cousinweb}
\textbf{\ta{தர்க்கம்}} (\emph{tarkkam}) is borrowed from Sanskrit \emph{tarka}, \textquotedblleft{}inquiry; conjecture\textquotedblright{} (CDIAL 5714), the name Indian thought gave to the study of reasoning.
\end{cousinweb}

\begin{grammarlens}[title={What you've built}]
One more word for giving reasons and opinions.
\end{grammarlens}

\subsection*{Your turn}
\begin{itemize}
\raggedright
  \item \emph{Say it:} \emph{tarkkam}
  \item \emph{Say it:} \emph{tarkkam}, once more
  \item \emph{Say it:} \emph{avaṉ coṉṉatil tarkkam illai.} — There is no logic in what he said.
  \item \emph{Recall:} say the Tamil for evidence, then the Tamil for to prove, then say \emph{tarkkam} again
\end{itemize}

\begin{tcolorbox}[breakable,colback=teal!4,colframe=teal!35!black,title={Before you move on}]
What does \ta{தர்க்கம்} mean? (\textquotedblleft{}Logic; reasoning\textquotedblright{}.) Say it once more.
\end{tcolorbox}
\section[\texorpdfstring{\ta{அடிப்படை} (aṭippaṭai)}{அடிப்படை (aṭippaṭai)}]{\ta{அடிப்படை} (aṭippaṭai) — a basis; a foundation; basic}

\label{lesson:TA-C347-atippatai}



\begin{quote}
Before the new one: say the Tamil for to prove, then the Tamil for logic.
\end{quote}

\subsection*{You'll want to know: \ta{அடிப்படை}}
\textbf{\ta{அடிப்படை}} — \emph{aṭippaṭai} — \textquotedblleft{}a basis; a foundation; basic\textquotedblright{}.

\begin{quote}
\ta{இந்த} \ta{முடிவின்} \ta{அடிப்படை} \ta{என்ன}?
\end{quote}

\begin{quote}
\emph{inta muṭiviṉ aṭippaṭai eṉṉa?}
\end{quote}

\begin{quote}
What is the basis of this decision?
\end{quote}

Before another noun it means basic: \ta{அடிப்படை} \ta{உரிமை}, a basic right.

\begin{cousinweb}
\textbf{\ta{அடிப்படை}} (\emph{aṭippaṭai}) is a native Tamil compound built from \emph{aṭi}, \textquotedblleft{}foot, base, bottom\textquotedblright{} (DEDR 72), and \emph{paṭai}, \textquotedblleft{}layer, stratum\textquotedblright{} (DEDR 3852): the bottom layer that the rest stands on.
\end{cousinweb}

\begin{grammarlens}[title={What you've built}]
One more word for giving reasons and opinions.
\end{grammarlens}

\subsection*{Your turn}
\begin{itemize}
\raggedright
  \item \emph{Say it:} \emph{aṭippaṭai}
  \item \emph{Say it:} \emph{aṭippaṭai}, once more
  \item \emph{Say it:} \emph{inta muṭiviṉ aṭippaṭai eṉṉa?} — What is the basis of this decision?
  \item \emph{Recall:} say the Tamil for to prove, then the Tamil for logic, then say \emph{aṭippaṭai} again
\end{itemize}

\begin{tcolorbox}[breakable,colback=teal!4,colframe=teal!35!black,title={Before you move on}]
What does \ta{அடிப்படை} mean? (\textquotedblleft{}A basis; a foundation; basic\textquotedblright{}.) Say it once more.
\end{tcolorbox}
\section[\texorpdfstring{\ta{விளக்கம்} (viḷakkam)}{விளக்கம் (viḷakkam)}]{\ta{விளக்கம்} (viḷakkam) — an explanation}

\label{lesson:TA-C347-vilakkam}



\begin{quote}
Before the new one: say the Tamil for logic, then the Tamil for a basis.
\end{quote}

\subsection*{You'll want to know: \ta{விளக்கம்}}
\textbf{\ta{விளக்கம்}} — \emph{viḷakkam} — \textquotedblleft{}an explanation\textquotedblright{}.

\begin{quote}
\ta{ஆசிரியர்} \ta{தெளிவான} \ta{விளக்கம்} \ta{கொடுத்தார்}.
\end{quote}

\begin{quote}
\emph{āciriyar teḷivāṉa viḷakkam koṭuttār.}
\end{quote}

\begin{quote}
The teacher gave a clear explanation.
\end{quote}

Its verb is \ta{விளக்கு}, to explain, the same word as \ta{விளக்கு}, a lamp: an explanation sheds light.

\begin{cousinweb}
\textbf{\ta{விளக்கம்}} (\emph{viḷakkam}) is native Dravidian: DEDR 5496 lists Tamil \emph{viḷakku}, \textquotedblleft{}to make bright, to explain\textquotedblright{}, and \emph{viḷakkam}, \textquotedblleft{}clearness, light\textquotedblright{}, in the family of \emph{veḷ}, \textquotedblleft{}white, bright\textquotedblright{}.
\end{cousinweb}

\begin{grammarlens}[title={What you've built}]
One more word for giving reasons and opinions.
\end{grammarlens}

\subsection*{Your turn}
\begin{itemize}
\raggedright
  \item \emph{Say it:} \emph{viḷakkam}
  \item \emph{Say it:} \emph{viḷakkam}, once more
  \item \emph{Say it:} \emph{āciriyar teḷivāṉa viḷakkam koṭuttār.} — The teacher gave a clear explanation.
  \item \emph{Recall:} say the Tamil for logic, then the Tamil for a basis, then say \emph{viḷakkam} again
\end{itemize}

\begin{tcolorbox}[breakable,colback=teal!4,colframe=teal!35!black,title={Before you move on}]
What does \ta{விளக்கம்} mean? (\textquotedblleft{}An explanation\textquotedblright{}.) Say it once more.
\end{tcolorbox}
